% =====================================================================
% TmFeO3 2DCS — main text (PRB)
% Long-form derivations: ../tmfeo3_foundation.tex; supplement: supp.tex
% =====================================================================
\documentclass[aps,prb,twocolumn,superscriptaddress,longbibliography]{revtex4-2}

\usepackage{amsmath,amssymb,bm}
\usepackage{graphicx}
\usepackage{xcolor}
\usepackage{hyperref}
\hypersetup{colorlinks=true,linkcolor=blue,citecolor=blue,urlcolor=blue}

\newcommand{\Fe}{\mathrm{Fe}}
\newcommand{\Tm}{\mathrm{Tm}}
\newcommand{\ket}[1]{|#1\rangle}
\newcommand{\bra}[1]{\langle #1|}
\newcommand{\ve}[1]{\boldsymbol{#1}}
\newcommand{\qafm}{q_{\rm AFM}}
\newcommand{\qfm}{q_{\rm FM}}
\newcommand{\Mnl}{M_{\rm NL}}

\begin{document}

\title{Two faces of the rare-earth--iron exchange:\\
two-dimensional terahertz spectroscopy of TmFeO$_3$ as vertex
spectroscopy of a dressed singlet ion}

\author{Working draft}
\affiliation{---}

\date{\today}

\begin{abstract}
One-dimensional spectroscopy measures the \emph{spectrum} of a
Hamiltonian; two-dimensional coherent spectroscopy (2DCS) measures its
\emph{vertices}---the individual nonlinear terms that convert one
coherence into another.  We show that in TmFeO$_3$ this correspondence
closes completely, because a single symmetry principle organises the
entire coupling space.  The Tm$^{3+}$ crystal-field levels are
non-Kramers singlets: in a real singlet basis time reversal makes every
angular-momentum matrix element imaginary, so the iron exchange field,
acting \emph{linearly}, reaches only the time-odd quadratures of the
crystal-field coherences---and that coupling fails on the data
threefold.  First-order exchange coupling to the crystal-field spectrum
is symmetry-blocked, and everything measured across the two subsystems
is a \emph{second-order dressing} with exactly two faces: the exchange
field squared, which modulates the crystal-field energy (a single
tensor $W$ that performs every coherence conversion but never
radiates), and the exchange field paired with the light field, which
builds an order-induced electric dipole
$g_d(\langle F_x\rangle{+}\delta F_x)\lambda^2$ (a single constant
that drives, radiates, and emits internal magnon--crystal-field sum-
and difference-frequency tones).  Driven by the measured pulse with
measured linewidths and dipole weights, these two objects plus the
bare subsystems and one self-absorption filter reproduce the
polarisation-resolved 2DCS atlas, with every peak assigned, every
structural absence a measurement, and every component validated by a
leave-one-out necessity audit.  Because $\partial^2E_{\rm
CEF}/\partial S\,\partial S$ is the same coupling surface whose
thermal average drives the spin-reorientation transition, the atlas
constitutes a coherent, mode-resolved, real-time measurement of the
reorientation mechanism itself, and the entire dressed sector must
collapse together at the transition.
\end{abstract}

\maketitle

\section{Introduction}

Linear response cannot distinguish two Hamiltonians with the same normal
modes: absorption measures pole positions and dipole weights, and any
coupling that merely hybridises modes is absorbed into the mode basis.
What it can never see is a \emph{conversion}: a coherence prepared on one
transition and re-emitted on another.  That is precisely what a cross
peak in 2DCS certifies~\cite{Lu2017,Mukamel}, and a conversion requires a
specific anharmonic vertex.  A cross peak at $(\omega_\tau,\omega_t)$ is
therefore not a feature to be fitted but a \emph{measurement of one
coupling term}: the delay axis names the stored coherence, the detection
axis the emitter, and the signed side constrains the pathway order.  If
symmetry makes the space of candidate vertices finite---and in a crystal
it does---a polarisation-resolved 2D atlas can be read as the Hamiltonian
itself.

The rare-earth orthoferrites are the natural target for this program.
Two qualitatively different subsystems share one THz octave: below the
spin reorientation the Fe sublattice of TmFeO$_3$ orders in $\Gamma_2$
$(F_x,C_y,G_z)$~\cite{Yamaguchi,Kimel2004,Zhang2016}, with
quasi-ferromagnetic and quasi-antiferromagnetic magnons at $\qfm=0.38$
and $\qafm=0.90$\,THz, while the three lowest Tm$^{3+}$ crystal-field
(CEF) singlets form a qutrit with $E_{12}=0.50$, $E_{23}=0.70$,
$E_{13}=1.20$\,THz.  The rare-earth--iron coupling in this family is
not a spectroscopic curiosity: it is the engine of the
spin-reorientation transition, in which the temperature-dependent
occupation of the rare-earth levels rotates the Fe anisotropy, and it
is the handle by which THz pulses have been used to steer the Fe order
coherently~\cite{Baierl2016,Schlauderer2019}.  Yet that coupling has
essentially always been treated thermodynamically---as a free-energy
surface---or as an empirical anisotropy modulation.  A single-cycle THz
pulse covers all five modes of both subsystems at once, and
measurements in the cross- and same-polarised channels of two field
geometries return distinct, reproducible peak inventories.  This paper
shows that the inventories can be read, term by term, as the coherent,
mode-resolved anatomy of the rare-earth--iron exchange---and that one
symmetry principle fixes what can appear.

\section{The dressing principle}
\label{sec:principle}

The organising fact is that the Tm$^{3+}$ CEF levels are nondegenerate
singlets of a non-Kramers ion.  Nondegenerate eigenstates of a real
crystal-field Hamiltonian may be chosen real; in that basis every
matrix element of the angular momentum $\ve J$ is purely imaginary and
antisymmetric.  Projected onto the qutrit, $\ve J$ therefore has
support \emph{only} on the three time-odd Gell-Mann coordinates
$\lambda^{2,5,7}$---the imaginary parts of the three coherences---and
the ion has rigorously no permanent moment.  Two consequences follow
immediately.

\emph{First-order exchange is blocked.}  The Fe sublattice exerts an
exchange field $\ve h_{\rm ex}(\ve S)$ on the ion; acting linearly it
produces the bilinear $K^-\,\ve S\cdot(\text{projected }\ve J)$, which
can only touch $\lambda^{2,5,7}$.  The data exclude this coupling
threefold: linear in $\ve S$ it hybridises modes but cannot convert
(no cross peak), its channels cancel over the four-sublattice orbit in
$\Gamma_1$ combinations to machine precision, and its one allowed
component reorients the ground state.  Its most direct fingerprint is
negative: a static mixer drags the dressed $E_{12}$ line with it, and
the measured 2D transfer row sits at the 1D absorption frequency
\emph{exactly}, at every tested coupling strength (with the bare
splitting recalibrated at each strength so the comparison is fair).
The singlet ion is protected from its magnetic environment at first
order---the same protection that makes non-Kramers ions poor qubits
makes them clean spectroscopic reporters.

\emph{Everything measured is second-order dressing, and there are
exactly two faces.}  If the environment cannot enter linearly, the
leading surviving objects are quadratic, and the qutrit's two
$1\!\leftrightarrow\!2$ quadratures each admit exactly one:
\begin{align}
  H_{W}&=\ve S^{T}W\,\ve S\;\lambda^1,
  &
  H_{g}&=-\,g_d\,E_z\,(F_x\,\lambda^2).
  \label{eq:faces}
\end{align}
The spin pair $SS$ is $\mathcal T$-even and lands on the
$\mathcal T$-even quadrature $\lambda^1$: this is the
\emph{energy face}---the exchange modulation of the crystal field,
$W\sim\partial^2 E_{\rm CEF}/\partial S\,\partial S$.  In $\Gamma_2$
symmetry leaves two components, $W^{yy}_1 S_y^2\lambda^1$ and
$W^{xz}_1 S_xS_z\lambda^1$, and the data are consistent with a single
scale $W^{yy}_1=W^{xz}_1=W$; the leave-one-out runs below show the two
components sharing every conversion task rather than owning separate
peaks---one tensor, not two mechanisms.  Having no field leg, the
energy face never radiates: it is the converter.  The light--spin pair
$E_zS$ is the \emph{moment face}: with the spin leg condensed on the
$\Gamma_2$ order parameter it is an order-induced electric dipole of
the $1\!\leftrightarrow\!2$ transition,
$d_z^{\rm eff}=g_d(\langle F_x\rangle+\delta F_x)\lambda^2$, reaching
the $\mathcal T$-odd quadrature through the $\mathcal T$-odd condensed
spin.  It is the one legal way an electric field talks to $\lambda^2$,
and one constant serves three ways: it \emph{drives} (the
$E\!\parallel\!c$ writing of the $E_{12}$ coherence), it
\emph{radiates linearly} (the $(\qafm,E_{12})$ readout), and its
retained fluctuation $g_d\,\delta F_x\lambda^2$ is a bilinear emitter
that radiates internal sum- and difference-frequency tones at
$\qafm\pm E_{12}$---magnon--crystal-field SFG and DFG, the
non-degenerate siblings of on-site second-harmonic generation.

This is the entire cross-subsystem sector.  Everything else in the
atlas is bare structure: the Fe magnon pair coupling (anisotropy plus
Dzyaloshinskii--Moriya, entirely within the Fe sector), the bare
driven qutrit with its measured dipoles and its on-site
hyperpolarisability $\beta\lambda^1\lambda^2$ (the SHG member of the
same quadratic-moment family, permitted because the Tm $4c$ site lacks
inversion; Fe at centrosymmetric $4b$ is excluded from the whole
family), and one propagation filter.  And the energy face is not a new
object in orthoferrite physics: $\partial^2E_{\rm CEF}/\partial
S\,\partial S$ is the coupling surface whose thermal average---through
the temperature-dependent rare-earth level occupations---rotates the
Fe anisotropy and drives the spin-reorientation
transition~\cite{Yamaguchi,Kimel2004,Schlauderer2019}.  The 2DCS atlas
measures the same surface coherently: mode-resolved, phase-resolved,
in both directions, on picosecond timescales.

\section{Model, and where it is exact}

The Fe$^{3+}$ ($S=5/2$) sublattice is classical, with the calibrated
orthoferrite Hamiltonian (exchange, single-ion anisotropy,
Dzyaloshinskii--Moriya) reproducing $\qfm$, $\qafm$ and $\Gamma_2$.  Each
Tm$^{3+}$ is truncated to its three lowest singlets and carried as the
Gell-Mann vector $\lambda^a=\langle\Lambda^a\rangle$.  Because every
single-ion term is linear in the $\Lambda^a$, the classical
$\vec\lambda$ equation of motion \emph{is} the exact von Neumann
equation of the driven, damped qutrit, and a Boltzmann ensemble over
level seeds is the exact thermal density matrix [SM, Sec.~S1].  The
Fe--Tm sector is built by projecting $\ve S^{T}\!A\,\ve J$ into the
qutrit in a transported gauge over the four Tm sublattices [SM,
Secs.~S2--S3]; the vertex catalogue and the exclusions behind
Eq.~(\ref{eq:faces}) are exhaustive over that space [SM, Sec.~S4].
All spectra use the measured pump (spectrum peaked at $0.53$\,THz, no
DC), measured linewidths as damping rates, and measured oscillator
strengths as emission weights.  Peaks are identified blindly as 2D
local maxima with prominence $\geq1.5$.

\section{Symmetry assigns each polarisation one block}
\label{sec:blocks}

The Tm $4c$ site carries the mirror $m$ ($z\to-z$), under which
$d_{x,y}$ are even and $d_z$ odd, while $m_z$ is even and $m_{x,y}$ odd.
Together with the time-reversal statement of Sec.~\ref{sec:principle}
this fixes the entire optical interface.  The projected moment matrix
---an output of the $PJP$ construction, not an input---has $\lambda^2$
carrying only $\mu_z$ and $\lambda^{5,7}$ only $\mu_{x,y}$, which
forces $\ket1,\ket2$ to share a parity and $\ket3$ to be opposite:
\begin{center}\small
\renewcommand{\arraystretch}{1.25}
\begin{tabular}{lccl}
\hline
transition & magnetic & electric & owned by\\
\hline
$1\!\leftrightarrow\!2$ & $m_z$ ($\lambda^2$) & $d_{x}$ ($\lambda^1$) &
  $(E\!\parallel\!a,H\!\parallel\!c)$\\
$1\!\leftrightarrow\!3$ & $m_{x}$ ($\lambda^5$) & $d_z$ ($\lambda^4$) &
  $(E\!\parallel\!c,H\!\parallel\!a)$\\
$2\!\leftrightarrow\!3$ & $m_{x}$ ($\lambda^7$) & $d_z$ ($\lambda^6$) &
  $(E\!\parallel\!c,H\!\parallel\!a)$\\
\hline
\end{tabular}
\end{center}
Mirror parity would additionally permit the mirror-even $d_x$ a
diagonal (permanent) part on $\lambda^{3,8}$; a permanent moment does
not radiate at any transition frequency, and the detection operators
exclude the population generators throughout (a measured rectified
line at $\omega_t\!\to\!0$ in an $(E\!\parallel\!a,H\!\parallel\!c)$
channel would falsify this exclusion).  An analyser selects a
polarisation \emph{state}, hence both the magnetic and the electric
dipole radiating into it; there are only \emph{two} detection
operators in the atlas, and the two channels sharing one are
constrained to share it exactly (Table~\ref{tab:channels}).
Reciprocity---one operator per polarisation, driving and radiating
alike---then implies that \emph{cross-polarised detection reads
exactly the block the pulse did not drive}.  This is the structural
reason cross-polarisation isolates conversion, and it carries two
independent checks: the geometry-B subalgebra closure, where the
$c$-directed drive touches only $\mu^z_{12}$ so that
$\lambda^{4\text{--}7}$ vanish identically (verified: exactly $0$,
against $\lambda^{1,2,3}\sim2.6\times10^{-2}$); and in geometry A
$F_z\equiv0$ to $10^{-15}$, so the Fe sublattice is radiatively silent
in that cross channel and anything seen there at a magnon frequency
must be Tm.

\begin{table}[b]
\caption{What each channel detects.  An analyser setting selects a
polarisation \emph{state}; there are therefore only two detection
operators.  $w_{E1}=0.54$ is the $E1/M1$ balance, fixed by the
geometry-B transfer (see text); no population generator
($\lambda^{3},\lambda^8$) appears in either operator, and the
$1\!\leftrightarrow\!3$ transition is fully dark
($\mu_{13}=0$ in both polarisations, $d_z(1{\to}3)\leq0.006$).}
\label{tab:channels}
\centering\small
\renewcommand{\arraystretch}{1.3}
\begin{tabular}{ll}
\hline
\multicolumn{2}{l}{$(E\!\parallel\!a,H\!\parallel\!c)$:\quad
  $F_z+\mu\lambda^2
   +w_{E1}\,\lambda^1
   +\beta\lambda^1\lambda^2$}\\
\hline
A cross & $F_z\equiv0$: Tm only.  $\lambda^{1,2}$ resonant at $E_{12}$\\
B same & $F_z$ alive but $50\times$ smaller than the Tm terms\\
\hline
\multicolumn{2}{l}{$(E\!\parallel\!c,H\!\parallel\!a)$:\quad
  $F_x+\mu_{x}\lambda^7
   +w_{E1}'\lambda^6
   +g_d(\langle F_x\rangle{+}\delta F_x)\lambda^2$}\\
\hline
A same & CEF $\lambda^{6}$; $d_z^{\rm eff}$ (static + dynamical);
  $F_x$ bounded\\
B cross & $F_x$; Tm part $\equiv0$ by closure\\
\hline
\end{tabular}
\end{table}

Because $\lambda^1$ and $\lambda^2$ are quadratures of the same
coherence, their magnitude spectra coincide to $10\%$: in
\emph{detection} the atlas is insensitive to the $E1/M1$ balance
(varying $w_{E1}$ over a factor of $50$ moves every geometry-A cross
amplitude by at most $0.04$ once $\beta$ is refitted).  It is the
\emph{drive} that fixes it: the electrically written coherence feeds
the geometry-B transfer linearly in $w_{E1}$, and the measured ratio
returns $w_{E1}=0.54$ (in units of $\mu$).  The construction is
strictly reciprocal at every level the data can test: operator
content and within-family dipole ratios are shared exactly between
drive and readout, as is the propagation filter below.  Three
cross-family scales remain as single calibrated numbers used
consistently on both sides: $w_{E1}$ (measured); $g_d$ (one parameter
with a drive avatar, a linear-emission avatar, and a
condensation-ratio correction $s=0.42$ on its bilinear piece, the
trajectory value overestimating it by about two---the model's
canting-angle uncertainty); and the Fe:Tm emission efficiency
$c_{\rm Fe}$, bounded from \emph{above} by the non-observation of the
$(\qafm,\qafm)$ magnon diagonal---a null measurement that is itself
the calibration.  The last ingredient is propagation: with $f=8.7$
the sample is optically thick at the $E_{13}$ line (depth $\simeq6$)
and mildly at $E_{12}$ (depth $\simeq1$), and the same transmission
filters the incoming field and the emitted $E1$ light---magnetic-dipole
radiation passes freely.  This one filter, applied identically both
ways, is itself reciprocal.

\begin{figure*}[t]
\includegraphics[width=\textwidth]{figs/atlas_final.png}
\caption{The atlas: four polarisation channels from one Hamiltonian at
one operating point.  Dashed guides mark the five mode energies;
crosses are blind-census maxima; linear amplitude colour scale; the
$\omega_\tau=0$ pump--probe row is retained in the geometry-A panels
(auto-scaled to the measured relative amplitude) and excluded from the
census.}
\label{fig:atlas}
\end{figure*}

\section{The pathways, peak by peak}

We now give every feature of Fig.~\ref{fig:atlas}: what makes it
possible, how the excitation reaches the emitter, and what would falsify
the assignment.  Throughout, the delay axis carries a coherence that
\emph{never radiates}---it need only be excitable and able to store
phase---while the detection axis is constrained by the emitting
multipole.  That asymmetry does most of the work.

\subsection{Geometry A, cross-polarised}

\emph{$(\qafm,E_{12})=1.00$: the energy face, read forward; measures
$W$.}  The $H\!\parallel\!a$ Zeeman field launches the $\qafm$
magnon, which precesses at $0.90$\,THz through the delay.  It cannot
radiate here---it has no $c$-axis moment---but it need not: it modulates
the Tm crystal field through $W^{yy}_1S_y^2\lambda^1$.  Each pulse
supplies one magnon leg, so the $\Mnl$ subtraction keeps the cross term,
which rectifies,
\begin{equation}
  h^1_{\rm NL}=W^{yy}_1A^2\big[\cos\qafm\tau-\cos\qafm(2t+\tau)\big]
  \theta(t),
\end{equation}
the difference piece being a step switched on at the probe that carries
the delay phase but no $t$ dependence.  Fed into the $E_{12}$ oscillator
it gives $\Mnl\propto\cos(\qafm\tau)[1-\cos E_{12}t]$: a factorised,
rank-one peak.  \emph{Why it is possible at all}: an E1-dark magnon is
made visible by borrowing the Tm dipole.  \emph{Where the energy comes
from}: $\delta S_y^2$ contains only DC and $2\qafm$, so between free
modes this vertex could never emit at $\omega_{12}$; the quantum is drawn
from the pulse edge, an intra-pulse difference-frequency process.  This
is verified rather than asserted---stretching the pulse from
$\sigma=0.25$ to $1.0$ at fixed amplitude collapses the peak by $4.5$
decades, while switching off the direct Tm drive changes it by
$<10^{-4}$.  \emph{Signature}: $\propto A_{\rm Fe}^2$, temperature-flat
below $\sim10$\,K, CEF linewidth along $\omega_t$.

\emph{$(E_{12},2E_{12})=0.96$: second-harmonic emission; measures
$\beta$.}  The pump stores $E_{12}$, setting the delay axis at $0.50$;
the detection axis is its second harmonic at $2E_{12}=1.0000$\,THz.  A
moment linear in the qutrit coordinates cannot radiate a harmonic;
the $4c$ site's lack of inversion permits the quadratic term
$m_z=\mu\lambda^2+\beta\lambda^1\lambda^2$, and Fe at $4b$ is excluded
from the mechanism.  Evaluated on the trajectories the quadratic
channel is remarkably clean, contributing essentially one peak,
$5.8\times$ its own $(E_{12},E_{12})$ diagonal, identical in local and
global sublattice frames.  The leave-one-out audit confirms the
assignment surgically: the harmonic survives removal of the Zeeman
drive and of both $W$ components at $1.00$--$1.03$, and dies to
$0.005$ without the electric $E_{12}$ writing.  \emph{Discriminators}:
the line sits at exactly $2E_{12}$ and \emph{tracks} $E_{12}$, so it
does not soften at the spin reorientation as any magnon-derived
feature would; it is one field order above the main peak; its width is
twice the $E_{12}$ linewidth.

\emph{$(\qafm,\qafm)=0.28$: the magnon diagonal, seen through the Tm
ion.}  Since $F_z\equiv0$ the magnon cannot radiate here at all.  What
radiates is the Tm dipole, forced at the magnon frequency by the
energy face: $W^{xz}_1$ contributes $G_z\,\delta S_x\,\lambda^1$,
linear in the magnon fluctuation, and $W^{yy}_1$ contributes the
pump-magnon$\times$probe-magnon cross term---the leave-one-out runs
apportion the peak between them ($W^{yy}_1$-dominant at the operating
point: removing it leaves $0.43$, removing $W^{xz}_1$ leaves
$0.64$).  Either way a driven oscillator responds at the driving
frequency, not its own, and the identification is verified in the
line itself: the sideband sits at $0.8999$\,THz---the magnon, not the
$\qfm+E_{12}=0.876$ combination---and the transfer coefficient
$|\lambda^2|_{0.90}/|m_x|_{0.90}=0.142$--$0.155$ is drive-independent
over a $20\times$ range, the hallmark of a forced linear response
evaluated off resonance.

\emph{$(\qafm,2\qafm)$: prediction.}  The same quadratic moment must
double \emph{any} frequency present in the Tm dipole, so the
magnon-forced component radiates its harmonic at $1.79$\,THz.  Its
observation confirms the hyperpolarisability independently of the
$2E_{12}$ peak; its absence with $(E_{12},2E_{12})$ present would
falsify the assignment.

\emph{$(-\qafm,E_{12})=0.50$ and $(\pm E_{12},E_{12})=0.32$.}  The first
is not a separate mechanism: the $W$ vertex is symmetric under the two
time orderings of the magnon coherence, so the signal appears on both
signed sides with the non-rephasing branch stronger.  The second is the
one-quantum diagonal, the directly driven $E_{12}$ read through the same
dipole; its size relative to the cross peaks is a sensitive fluence
diagnostic, falling from $1.45$ at $A_{\rm Fe}=1.2$ to $0.16$ at $0.12$,
which is how the operating point was fixed.

\subsection{Geometry B, cross-polarised}

Here the analysed moment is $m_x$, which the $\qafm$ magnon carries
directly, so both peaks emit at $0.90$\,THz through the Fe magnetic
dipole and the Tm ion is radiatively absent (the closure above).

\emph{$(\qfm,\qafm)=1.00$: magnon--magnon; measures the Fe sector
alone.}  $H\!\parallel\!c$ launches the quasi-ferromagnetic magnon,
which precesses at $0.38$ through the delay; single-ion anisotropy and
the Dzyaloshinskii--Moriya interaction couple the branches, so the stored
$\qfm$ phase is re-emitted on $\qafm$.  No Fe--Tm coupling enters at any
vertex, which is why this peak is temperature-flat: a pure two-magnon
process, independent of the qutrit populations.

\emph{$(E_{12},\qafm)=0.42$: the energy face, read backward;
measures $w_{E1}W$.}  $E\!\parallel\!a$ owns the
$1\!\leftrightarrow\!2$ block, so the pulse's \emph{electric} field
writes the $E_{12}$ coherence through $d_x\lambda^1$; the energy face
then hands the stored coherence to the quasi-antiferromagnetic magnon,
which radiates through the Fe $m_x$ dipole.  No field enters the
conversion: the polarisation ownership of $d_x$ performs the geometry
selection a field gate would otherwise be needed for---this drive
component does not exist in geometry A, which is why the conversion
never contaminates the clean $W^{yy}_1$ peak there.  Three facts pin
the assignment.  First, the \emph{quadrature} matters: writing the
same coherence magnetically (through $\mu_z\lambda^2$) converts far
less cleanly, flooding the map with drive-squared diagonals, whereas
the electric writing at the master $W=0.01$ adds the transfer row at
the \emph{unshifted} line position with strays below $0.17$.  Second,
the interference of the two drive quadratures places the row: with
$\mathrm{sign}(d_x\mu_z)<0$ it sits at $\omega_\tau=0.52$, exactly
where it is measured---the relative sign is thereby a prediction for
the crystal-field wavefunctions.  Third, the electric component's
amplitude is linear in $w_{E1}$, and matching the measured
$0.43$--$0.46$ returns $w_{E1}=0.54$.  The leave-one-out runs quantify
the pathway shares: removing the $d_x$ writing halves the transfer and
removing $W^{xz}_1$ costs a quarter, the remainder flowing through the
magnetically written quadrature and the $W^{yy}_1$ mixing route
($\qfm{+}E_{12}=0.88$, unresolved from $\qafm$; either route
ultimately rings the real $\qafm$ mode, and the emission stays pinned
at $0.900$ in every ablation)---the electric leg is the largest single
element, and it is the $w_{E1}$-bearing one.
\emph{Signature}: falls as the population difference $n_1-n_2$
($0.46\to0.45\to0.40\to0.36$ at $0/5/8/10$\,K), the sharpest thermal
discriminator in the atlas against the temperature-flat magnon--magnon
peak; and, riding the order parameter through the static vertex
($\propto G_z$), the transfer must collapse at the spin reorientation.

\emph{The post-drive buildup.}  The measured geometry-B signal continues
to grow after the pulse.  Within the present model no channel reproduces
growth at the peak's own scale.  Every coherent pathway inherits its
source's lifetime---because $B\!\parallel\!c$ never drives $\qafm$
directly, all geometry-B magnon amplitude is converted, and with
$T_2(E_{12})=8$\,ps everything coherent is over within $\sim10$\,ps.
The one reservoir that outlives every $T_2$---the CEF populations, fed
incoherently by the dissipated pulse energy and acting on the magnon
through the population channel $W_3$, the mechanism of thermally driven
spin reorientation---does produce a persistent, drive-uncorrelated tail
in the solver, but roughly four decades below the impulsive peak.  The
buildup therefore stands as the one timing observable outside this
coupling space: a confirmed growth of the transfer peak's own
amplitude would require a reservoir this Hamiltonian does not contain.

\subsection{Geometry A, same-polarised}

The analysed moment is $m_x$, and three contributions reach it, each
required by symmetry: the Tm CEF moment (the $d_z$ line at $E_{23}$,
which also carries the cascade tones below; the $1\!\to\!3$ member is
bounded at $d_z\leq0.006$ by the weak $1.15$ wing and dropped from the
minimal operator), the order-induced dipole
$g_d(\langle F_x\rangle+\delta F_x)\lambda^2$, and the Fe $F_x$
itself.  The weights are fitted to the digitised measurement with the
cross channels held fixed; the Fe emission efficiency $c_{\rm Fe}$ is
\emph{bounded from above by a null}, the absence of the magnon
diagonal.  The result is Table~\ref{tab:samepol}.

\begin{table}[b]
\caption{Geometry-A same-polarised channel: measured amplitudes against
the final model (overlap-gated delay apodisation; self-absorption of
the $E1$ emission at the $E_{13}$ line, optical depth $\simeq6$ on both
the emission column and the delay-axis FID, and depth $\simeq1$ at
$E_{12}$).  The $\mu^z_{13}$ \emph{and} $\mu^x_{13}$ weights fit to
zero (fully dark $1\!\to\!3$).}
\label{tab:samepol}
\centering\small
\renewcommand{\arraystretch}{1.25}
\begin{tabular}{lccl}
\hline
peak & exp. & model & pathway\\
\hline
$(E_{12},0)$ & $1.00$ & --- & rectified storage (see text)\\
$(E_{12},E_{23})$ & $1.00$ & $1.00$ & coherence$\to$pop.$\to\rho_{23}$\\
$(E_{12},0.40)$ & $0.45$ & $0.40$ & cascade tone $2(E_{23}{-}E_{12})$, adjacent to $\qfm$\\
$(E_{12},0.90)$ & $0.40$ & $0.43$ & cascade tone $2E_{23}{-}E_{12}$, degenerate with $\qafm$\\
$(0,E_{12})$ & $0.40$ & --- & pump--probe row\\
$(\qafm,E_{12})$ & $0.30$ & $0.39$ & magnon$\to$Tm via $d_z^{\rm eff}$\\
$(\qafm,\qafm)$ & absent & $0.13$ & magnon diagonal; absence bounds $c_{\rm Fe}$\\
$(E_{13},E_{23})$ & $0$ & $0.00$ & 2Q row, killed by FID reabsorption\\
$(E_{12},1.15)$ & $0.15$ & $0.02$ & $E_{13}$ wing; bounds $d_z(1\to3)\leq0.006$\\
$(\qafm,1.41)$ & $0.30$ & $0.31$ & dynamical $d_z^{\rm eff}$: SFG tone $\qafm{+}E_{12}$\\
$(\qafm,0.40)$ & obs. & $0.43$ & its DFG tone $\qafm{-}E_{12}$\\
$(E_{12},1.41)$ & obs. & $0.36$ & its $E_{12}$-delay partner\\
$(-E_{12},E_{23})$ & $<0.15$ & $0.7$ & \textbf{open} (residual R1, point-ion floor)\\
\hline
\end{tabular}
\end{table}

\emph{The fully dark $1\to3$ dipole---measured, not assumed.}  The
measured map contains delay frequencies $\omega_\tau\in\{E_{12},\qafm\}$
only.  A bright $\mu^z_{13}$ would store an $E_{13}$ coherence at first
order and produce an $(E_{13},E_{23})$ peak with \emph{exactly} the same
dipole product as the observed $(E_{12},E_{23})$: the absence of the
whole row is a selection rule.  The data force more: the analyser for
this channel physically contains $\mu_x\lambda^5$, and at the
crystal-field value $\mu^x_{13}=2.39$ that term would exceed the entire
measured map by two orders of magnitude---so $\mu^x_{13}$, too, fits to
zero.  The $1\!\to\!3$ transition is magnetically dark in \emph{every}
polarisation; its large measured oscillator strength ($f=8.7$) is
electric, and it is precisely this strong $E1$ line that makes the
sample opaque at $1.20$\,THz (below).

\emph{$(E_{12},E_{23})=1.00$: the dominant transfer, via a population.}
With the single-action route closed by the dark dipole, the surviving
pathway acts twice: one probe interaction converts the stored $E_{12}$
coherence into a level-2 \emph{population} that still carries the
$\omega_\tau=E_{12}$ phase, and a second converts that population into
the emitting $E_{23}$ coherence.  Being co-rotating it lands
non-rephasing, the observed side, and it involves $\mu_{13}$ at no
vertex---which is why it survives while its $E_{13}$ partner is
throttled.

\emph{$(E_{12},0)=1.00$: the rectified line.}  Taken at its
zero-frequency output the same double action rectifies the stored
coherence into a slow net magnetisation of the detection window.  This
is the strongest measured feature; in the 2D maps it is suppressed by
construction, since the $t\geq3$ window with Gaussian apodisation removes
$\omega_t\to0$, so it is analysed separately through
$S_{\rm DC}(\tau)=\langle\Mnl\rangle_t$, where it appears as a pair of
sharp lines at $\pm0.49$\,THz---symmetric in the delay sign, as a
population grating must be ($\cos E_{12}\tau$), and as the measurement
finds ($1.00/0.85$).  Its \emph{magnitude} is a physical statement of
the first importance: under coherent pumping alone the model's
rectified moment is orders of magnitude below the resonant peaks,
whereas the measurement puts it at the top of the map.  The line is
therefore the optical readout of the \emph{incoherent} population
channel---absorbed pulse energy returning through the CEF populations
and $W_3$ to a quasistatic change of the ordered moment---the same
reservoir the geometry-B post-pulse growth demands, and its strength
relative to $(E_{12},E_{23})$ calibrates $c_{\rm heat}W_3$ directly.

\emph{$(E_{12},0.40)=0.45$ and $(E_{12},0.90)=0.40$: cascade tones of
the driven qutrit, masquerading at the magnon frequencies.}  The
necessity audit (removing the Fe Zeeman leg entirely and re-running)
leaves both features at $0.89$--$0.97$ of their baseline: they are
not magnon features at all but intrinsic tones of the three-level
cascade, $2(E_{23}{-}E_{12})=0.40$ and $2E_{23}{-}E_{12}=0.90$\,THz,
radiated on the $\lambda^6$ ($d_z$) line and numerically coincident
with $\qfm=0.38$ and $\qafm=0.90$.  Both die without the $\lambda^7$
drive leg ($0.05$--$0.08$), confirming the double-$\rho_{23}$
pathway; the genuine Tm$\to$magnon transfer is the
$\lesssim\!10\%$ Zeeman-sensitive remainder.  The measurement can
separate the hypotheses directly: the low tone sits at $0.40$, more
than a linewidth above $\qfm=0.38$, and a cascade tone survives the
spin reorientation with CEF width while a magnon sideband would
soften and sharpen.  The emission is in either case electric and
rare-earth, insensitive to the Fe weight ($0.46\to0.43$ when
$c_{\rm Fe}$ is reduced fourfold).

\emph{$(\qafm,E_{12})=0.30$: the moment face, read linearly.}  The
magnon occupies the delay axis and the Tm ion radiates at $E_{12}$ into
the \emph{same}-polarised channel.  That emission is possible only
through $d_z^{\rm eff}$---and reciprocity demands it, because the same
$d_z^{\rm eff}$ is what lets $E\!\parallel\!c$ create the $E_{12}$
coherence in the first place.  We had earlier reported this partner as
absent and treated the mismatch as an open tension; the digitised map
shows it present at $0.30$ of maximum, and the model reproduces it at
$0.39$, fed jointly by the condensed readout and the low-frequency
side of the SFG/DFG tones.  This peak is the direct measurement of
$g_d\langle F_x\rangle$.

\emph{The SFG/DFG triplet: $(\qafm,1.41)$, $(\qafm,0.40)$,
$(E_{12},1.41)$---the moment face, read bilinearly.}  Condensing the
spin leg of $H_g$ on $\langle F_x\rangle$ was a truncation: the full
emitting moment is $g_d(\langle F_x\rangle+\delta F_x)\lambda^2$, and
the fluctuation piece is a product emitter---the stored $\qafm$ magnon
beating against the $E_{12}$ coherence on one dipole---radiating
internal sum- and difference-frequency tones at
$\qafm\pm E_{12}=1.40/0.40$\,THz with the magnon on the delay axis,
plus an $E_{12}$-delay partner at $1.40$.  All three tones are
measured.  This channel is singular in the entire coupling space: of
some forty conversion vertices, emission moments and quadratic
products catalogued, it is the \emph{only} one whose map is led by the
$\qafm$-delay row rather than the $E_{12}$-delay row, because both of
its factors are written at first order by the same pulse---the stored
magnon pays no conversion penalty, and the audit confirms it (the
tones survive removal of both $W$ components at $0.95$--$1.02$ and die
only with the Zeeman or electric drive legs).  One number is fitted
(the condensation-ratio correction $s=0.42$), after which the sum tone
lands at the measured amplitude, the difference tone merges into the
measured $(\qafm,E_{12})$ blob, and the partner appears at $0.36$.
The sum tone escapes one absorption linewidth \emph{outside} the
opaque $E_{13}$ window (transmission $0.68$ at $1.39$ against $0.002$
at line centre), which is why it is visible while the true $E_{13}$
line centre is dark.  Together with the on-site $\beta$ these tones
complete the $\chi^{(2)}$-like triad of the inversion-free rare-earth
site, \emph{in emission}: SHG ($\beta\lambda^1\lambda^2$, degenerate),
SFG and DFG ($g_d\,\delta F_x\lambda^2$, non-degenerate,
cross-subsystem)---the mixing element being a bilinear term of the
material's transition moment, not a crystal-optical $\chi^{(2)}$
acting on light.

\emph{$(\qafm,\qafm)$: absent---and the absence is a measurement.}
The magnon diagonal is the one feature of this channel that genuinely
requires Fe emission: a $\qafm$ delay coherence re-emitting at its own
frequency through the $F_x$ magnetic dipole, with no Tm conversion
involved.  The measured map does not contain it.  Since every other
magnon-frequency feature is Tm-carried, this null bounds the Fe
emission efficiency from above ($c_{\rm Fe}$ set at the bound puts the
diagonal at $0.13$, below the digitisation threshold, and moves
nothing else).  The physical statement is uniform across the atlas:
the Fe sublattice is loud in geometry B, where the magnon owns the
analysed $m_x$ directly and is driven resonantly, and quiet in
geometry A, where everything measured arrives through the rare-earth
ion ($F_z\equiv0$ makes the cross channel's Fe silence exact; the
small $c_{\rm Fe}$ makes the same-polarised channel's approximate).

\emph{The $E_{13}$ column: self-reversal.}  The two measured
$E_{13}$-emission features sit at $\omega_t=1.15$ and $1.25$\,THz---one
absorption linewidth \emph{below and above} the $1.20$ line---which is
the textbook signature of self-reversal: with $f=8.7$ the sample is
optically thick at $E_{13}$, the line centre of any $E1$ emission is
reabsorbed, and only the wings escape.  The same reabsorption acts on
the delay-axis FID, which is what removes the two-quantum
$(E_{13},E_{23})$ row \emph{exactly} (model $0.33\to0.00$ once the
filter is applied)---a row the free-ion model robustly predicts and the
measurement robustly lacks.

\subsection{Geometry B, same-polarised: a one-line prediction}

This channel is analysed with the \emph{same} operator as geometry-A
cross, so it is a prediction with no freedom left.  With the population
generators excluded from detection, three terms survive: $F_z$ (alive
here, but $\sim\!50\times$ below the Tm terms), the coherence readouts
$\mu\lambda^2$ and $w_{E1}\lambda^1$, and the hyperpolarisability
$\beta\lambda^1\lambda^2$.  The linear readouts are suppressed in
$\Mnl$---they carry the directly driven coherence, which the
subtraction removes at first order---so the channel collapses onto its
quadratic term.  The prediction is a \emph{single} broad feature: the
second harmonic $(E_{12},2E_{12})=1.00$, with everything else below
$0.07$---magnon features at the $10^{-2}$ level of $F_z$, and no
$E_{13}$ or $E_{23}$ emission by the subalgebra closure
$\lambda^{4\text{--}7}\equiv0$.  Because $\beta$ is already fixed by
the geometry-A cross parity, the channel is a zero-parameter test of
the hyperpolarisability: a map led by anything other than the
$2E_{12}$ ridge falsifies the quadratic moment, and an observed
rectified pair at $(\pm E_{12},\omega_t\!\to\!0)$ would falsify the
exclusion of the population readout from detection.

\section{The necessity audit}
\label{sec:audit}

Every component of the final model was subjected to a leave-one-out
test: drive legs and Hamiltonian vertices by re-running the solver
with one term removed, detection terms and propagation filters by
dropping one term of the composite.  Table~\ref{tab:audit} summarises.
The result is minimal in the audited sense: two cross-subsystem
objects ($W$, $g_d$), one on-site anharmonicity ($\beta$), one
measured drive ratio ($w_{E1}$) and one propagation filter carry the
entire coherent atlas.  Every other symmetry-allowed object is zero,
bounded, or inert: the field-assisted family is \emph{bitwise} inert
in these geometries (removing $\kappa^B$ reproduces the trajectories
to the last digit---its gating field component is never driven), the
population channel $W_3$ moves no coherent feature by more than
$5\%$ and enters only the incoherent observables, and the four
remaining quadratic moments of the $4c$ site are excluded by the
smallness of the measured $(E_{12},1.20)$ feature, since each would
radiate the stored $E_{12}$ coherence at the $E_{12}{+}E_{23}$ sum
frequency at line centre.

\begin{table}[t]
\caption{The leave-one-out audit (single-seed ablation runs; ratios of
feature amplitudes to the baseline).}
\label{tab:audit}
\centering\small
\renewcommand{\arraystretch}{1.2}
\begin{tabular}{lll}
\hline
removed & what dies (ratio) & verdict\\
\hline
Fe Zeeman leg & A-cross anchor $0.02$; tones $0.03$ & required\\
$d_z^{\rm eff}$ drive & anchors, SHG $0.005$; tones $0.1$ & required\\
$\mu_x\lambda^7$ drive & anchor $0.02$; cascade tones $0.05$ & required\\
$W^{yy}_1$ & A-cross anchor $0.03$ & required\\
$W^{xz}_1$ & sideband $0.64$; B transfer $0.76$ & required (shared)\\
$\beta$ (detect) & SHG $0.95\to0.01$ & required\\
$\lambda^6$ (detect) & same-pol anchor $0.06$ & required\\
$\delta F_x\lambda^2$ (detect) & SFG/DFG triplet $0.02$--$0.05$ & required\\
$T_{13}$ (column / row) & $0.67$ line-centre peak / 2Q row $0.32$ & required\\
$d_x$ drive (geom.\ B) & transfer $0.49$ & dominant leg\\
$\kappa^B$ & nothing (bitwise) & inert, exact\\
$W_3$ + thermal & coherent features $\leq5\%$ & atlas-inert\\
$\lambda^4$, $F_z$ (detect) & only the $1.15$ wing / nothing & bounded / spectator\\
\hline
\end{tabular}
\end{table}

The audit also corrected two assignments (already incorporated above):
the same-polarised $E_{12}$-row features at $0.40$ and $0.90$ are
driven-qutrit cascade tones, not Tm$\to$magnon transfers, and the
A-cross magnon sideband and geometry-B transfer are shared between the
two $W$ components rather than owned by one.  Both corrections
sharpen, rather than weaken, the dressing principle: the conversion
sector is one tensor, and the features that need no conversion at all
are exactly the ones that survive every ablation of it.

\section{What this says about rare-earth orthoferrites}
\label{sec:physics}

Three statements generalise beyond TmFeO$_3$.

\emph{The rare-earth ion is a coherent transducer, not a passive
thermal bath.}  In the standard picture of the orthoferrites the
$R$-ion enters through its free energy: thermally occupied CEF levels
shift the Fe anisotropy, and reorientation follows.  The atlas shows
the same coupling surface operating \emph{coherently}: the energy face
$W\sim\partial^2E_{\rm CEF}/\partial S\,\partial S$ converts quanta
between magnon and CEF coherence within picoseconds, in both
directions (the geometry-A anchor forward, the geometry-B transfer
backward), phase-resolved and mode-resolved.  The couplings that
reorient the statics and the vertices that convert the dynamics are
one object read at two temperatures of description---which is
presumably why THz control of the Fe order through the rare-earth
channel is possible at all~\cite{Baierl2016,Schlauderer2019}.  The
sharpest consequence is collective: every dressed feature in the
atlas---the transfer peaks, the $(\qafm,E_{12})$ readout, the SFG/DFG
triplet---rides $\langle F_x\rangle$ or $G_z$ and must collapse
\emph{together} at the reorientation, while the bare-structure
features (the magnon--magnon peak, the SHG line, the cascade tones)
survive it.  A temperature scan through $T_{\rm SR}$ cleaves the atlas
into its two halves along exactly the line the dressing principle
draws.

\emph{Kramers parity controls the atlas.}  The cleanliness of
TmFeO$_3$---sharp unshifted lines, perturbative assignable vertices,
a readable atlas---is the singlet protection of
Sec.~\ref{sec:principle} at work: first-order exchange is
symmetry-blocked, so nothing hybridises strongly and everything
converts weakly.  In a Kramers member of the family (Er, Yb) no such
protection exists: linear exchange coupling is allowed, the $R$ modes
and magnons form mixed magnon--CEF excitations, and the corresponding
2DCS atlas should be qualitatively different---dominated by
hybridisation shifts and shared-linewidth mixed modes rather than by
weak, assignable conversion vertices.  The non-Kramers orthoferrites
(Tm, Pr, Tb) are the clean laboratories; comparing a Kramers and a
non-Kramers member under the same protocol would turn this prediction
into a family-wide statement.

\emph{Order-induced moments carry mandatory fluctuation sidebands.}
The moment face is an order-induced dipole, and the audit forced its
un-truncation: if a transition moment exists only because an order
parameter condensed ($d_z^{\rm eff}\propto\langle F_x\rangle$), then
the same coupling necessarily radiates sum- and difference-frequency
tones with the order parameter's own fluctuations
($g_d\,\delta F_x\lambda^2$ at $\qafm\pm E_{12}$), with amplitudes
locked to the condensed readout by a single ratio.  This is a
universal signature, not a TmFeO$_3$ accident: any order-induced line
in any magnetoelectric---multiferroic electromagnons, exchange-induced
dipoles, field-induced moments---must carry these sidebands, and 2DCS
is the instrument that separates them from cascades, because the
delay-axis label distinguishes a stored order-parameter fluctuation
from a stored coherence.  Here the sidebands were not decoration: they
resolved the last unassigned peak in the spectrum.

\section{Falsifiable content}

The model constrains itself as follows.  Any peak emitting at $\qafm$
through the Fe dipole inherits the $15\times$ sharper magnon
linewidth; $(\qfm,\qafm)$ is temperature-flat while $(E_{12},\qafm)$
falls as $n_1-n_2$; the transfer row sits at the frequency of the 1D
absorption line exactly (a static mixer would place it below, at any
coupling strength); $\mathrm{sign}(d_x\mu_z)<0$, checkable against the
crystal-field wavefunctions; rotating either polarisation away from
$c$ must revive the $\omega_\tau=E_{13}$ rows; the $2E_{12}$ line
tracks $E_{12}$ and does \emph{not} soften at the spin reorientation,
whereas the entire dressed sector---the $d_z^{\rm eff}$ family
\emph{and} the $(E_{12},\qafm)$ transfer---must vanish there; the
geometry-B same-polarised map must be led by the $2E_{12}$ harmonic
essentially alone; and adiabatically long pulses collapse the
geometry-A cross peak.  Propagation adds its own fingerprints: a
lineout across $\omega_t=1.1$--$1.3$ at $\omega_\tau=\qafm$ must show
a self-reversal \emph{doublet} around a dark centre at $1.20$; the
geometry-B cross channel must carry a $d_z^{\rm eff}$ doublet at
$\omega_t=0.45/0.54$ at the $\sim\!0.17$ level; and every wing
amplitude scales with sample thickness, so a thin sample revives the
$E_{13}$ line centre and the two-quantum row together.  The emission
assignments carry lineshape and position fingerprints: the
same-polarised $E_{12}$-row tones must sit at $2(E_{23}{-}E_{12})=0.40$
and $2E_{23}{-}E_{12}=0.90$ with CEF width and survive the
reorientation---at the low tone a full linewidth above $\qfm=0.38$, a
direct positional test---and any recovered $(\qafm,\qafm)$ diagonal
would raise the $c_{\rm Fe}$ bound and must appear with the magnon
width.  The SFG/DFG triplet is over-constrained: all three tones come
from one product, so their ratios are locked, the sum tone must sit at
$\qafm{+}E_{12}$ exactly (at $1.20$ it would be a quadratic moment
instead, and the dark $E_{13}$ window separates the two), its position
must track the magnon as $\qafm$ softens, and the whole triplet
collapses at the reorientation with everything else that rides
$g_d\propto\langle F_x\rangle$.  The temperature scan through the
reorientation remains the decisive experiment: it cleaves the atlas
into dressed and bare halves and measures $g_d$ by the size of the
collapse.

One residual survives, named and characterised.  \emph{R1, the
rephasing remnant} $(-E_{12},E_{23})$: the point-ion qutrit has an
exact floor of $0.47$--$0.58$ of the dominant transfer (verified
against a standalone three-level integrator; no coupling in the
exhausted vertex space and no drive quadrature or sign removes it),
while the measured negative-delay side carries only the rectified and
$\qfm$ features---and the SFG/DFG tones add their own rephasing
partners there (model $0.5$--$0.7$), sharpening the same question.
This residual, the self-reversal wing amplitudes, and the geometry-B
$d_z^{\rm eff}$ doublet strength funnel into a single measurable
object: the sample's polarisation-resolved optical depths at the three
CEF lines.  One transmission measurement closes it or breaks the
propagation layer cleanly.

\section{Conclusion}

A single symmetry principle---the time-reversal protection of a
non-Kramers singlet ion---reduces the rare-earth--iron coupling space
of TmFeO$_3$ to two second-order dressings, one per quadrature of the
crystal-field coherence: an energy face
$W\,SS\,\lambda^1$ that converts and never radiates, and a moment face
$g_d\,E_z(F_x\lambda^2)$ that drives and radiates, linearly through
its condensed part and bilinearly through its fluctuations (internal
magnon--crystal-field SFG/DFG).  These two objects, the bare
subsystems, one on-site hyperpolarisability and one self-absorption
filter reproduce the measured polarisation-resolved 2DCS atlas from
measured inputs, with no field-assisted coupling and no free lineshape
parameters, and every component certified necessary by a leave-one-out
audit.  The absences do quantitative work alongside the peaks: the
missing $E_{13}$ delay row measures the fully dark $1\!\to\!3$ dipole,
the missing magnon diagonal bounds the Fe emission weight, and the
missing $(E_{12},1.20)$ sum-frequency partner excludes every remaining
quadratic moment of the rare-earth site.  Because the energy face is
the spin-reorientation coupling surface itself, the atlas is a
coherent, mode-resolved measurement of the mechanism that moves the
statics of this family---and its sharpest prediction is collective:
at the reorientation the dressed half of the atlas dies together,
while the bare half survives.  The program---bound the vertex space by
symmetry, assign every feature and every absence, audit every
component, and let the atlas overdetermine the model---applies to any
magnet whose 2D map spans distinct subsystems, and turns 2DCS from a
spectroscopy of modes into a spectroscopy of Hamiltonian terms.

\begin{acknowledgments}
Derivations, numerical protocol and the complete falsification ledger are
in the Supplemental Material; the long-form note and all run recipes
accompany the data.
\end{acknowledgments}

\begin{thebibliography}{9}
\bibitem{Lu2017} J. Lu, X. Li, H. Y. Hwang, B. K. Ofori-Okai,
  T. Kurihara, T. Suemoto, and K. A. Nelson,
  Phys.\ Rev.\ Lett.\ \textbf{118}, 207204 (2017).
\bibitem{Mukamel} S. Mukamel, \emph{Principles of Nonlinear Optical
  Spectroscopy} (Oxford University Press, New York, 1995).
\bibitem{Yamaguchi} T. Yamaguchi,
  J.\ Phys.\ Chem.\ Solids \textbf{35}, 479 (1974).
\bibitem{Kimel2004} A. V. Kimel, A. Kirilyuk, A. Tsvetkov, R. V.
  Pisarev, and Th. Rasing, Nature \textbf{429}, 850 (2004).
\bibitem{Zhang2016} \emph{Resolving the spin reorientation and
  crystal-field transitions in TmFeO$_3$ with terahertz transient},
  Sci.\ Rep.\ \textbf{6}, 23648 (2016).
\bibitem{MP2018} \emph{Terahertz magnon-polaritons in TmFeO$_3$},
  ACS Photonics \textbf{5}, 1375 (2018).
\bibitem{Baierl2016} S. Baierl \emph{et al.}, \emph{Nonlinear spin
  control by terahertz-driven anisotropy fields},
  Nat.\ Photonics \textbf{10}, 715 (2016).
\bibitem{Schlauderer2019} S. Schlauderer \emph{et al.}, \emph{Temporal
  and spectral fingerprints of ultrafast all-coherent spin switching},
  Nature \textbf{569}, 383 (2019).
\end{thebibliography}

\end{document}
